\documentclass[conference]{IEEEtran}

\usepackage{cite}
\usepackage{amsmath,amssymb}
\usepackage{graphicx}
\usepackage{booktabs}
\usepackage{siunitx}
\usepackage{xcolor}
\usepackage[hidelinks]{hyperref}

\graphicspath{
  {figures/phase_02/}
  {figures/phase_03/}
  {figures/phase_04/}
}

% Reusable figure command for ready and planned paper assets.
% Include this file after loading the graphicx package.
%
% Arguments:
%   #1 relative figure path
%   #2 image width, such as 0.95\linewidth
%   #3 placeholder description
%   #4 final figure caption
%   #5 semantic LaTeX label

\newcommand{\ProjectFigure}[5]{%
  \begin{figure}[tbp]
    \centering
    \IfFileExists{#1}{%
      \includegraphics[width=#2]{#1}%
    }{%
      \fbox{%
        \begin{minipage}[c][0.24\textheight][c]{0.90\linewidth}
          \centering
          \textbf{FIGURE PLACEHOLDER}\\[0.8em]
          \texttt{\detokenize{#1}}\\[0.8em]
          #3
        \end{minipage}%
      }%
    }
    \caption{#4}
    \label{#5}
  \end{figure}%
}


% Visible notes make incomplete prose difficult to mistake for final content.
\newcommand{\DraftingNote}[1]{%
  \par\noindent
  \fcolorbox{red!70!black}{red!5}{%
    \parbox{0.94\columnwidth}{\textbf{DRAFTING NOTE:} #1}%
  }\par
}

\title{Noise-Robust Urban Sound Classification: Evaluating Background Noise
and Audio Augmentation Across Deep Learning Models}

\author{
  \IEEEauthorblockN{Author One, Author Two, and Author Three}
  \IEEEauthorblockA{
    Department of Computer Science and Engineering\\
    University Name, City, Country\\
    \{author1, author2, author3\}@example.edu
  }
}

\begin{document}

\maketitle

\begin{abstract}
\DraftingNote{Write the abstract last. It should state the problem, UrbanSound8K
split, six trained configurations, four evaluation conditions, strongest
quantitative result, limitations, and conclusion without adding unsupported
claims.}
\end{abstract}

\begin{IEEEkeywords}
environmental sound classification, audio augmentation, background noise,
robustness, Log-Mel spectrogram, UrbanSound8K
\end{IEEEkeywords}

\section{Introduction}

\DraftingNote{Motivate environmental sound classification under real background
noise. Introduce the gap between clean test accuracy and controlled robustness,
then summarize the study contributions. Add only verified primary citations.}

This study investigates whether training-time audio augmentation reduces
performance degradation when urban sound classifiers are evaluated with
controlled background noise. Three architectures are compared: a compact
two-dimensional convolutional neural network (CNN), a convolutional recurrent
neural network (CRNN) with a bidirectional gated recurrent unit (BiGRU), and a
non-pretrained ResNet18 adapted to one-channel Log-Mel spectrograms.

\subsection{Research Questions}

\begin{enumerate}
  \item Does training-time audio augmentation reduce the loss of classification
  performance when test audio is corrupted by background noise?
  \item Is the augmentation effect consistent across CNN, CRNN, and ResNet18
  architectures?
\end{enumerate}

\subsection{Contributions}

\DraftingNote{Refine this list after the rest of the paper is stable. Keep the
claims within the one-fold, one-training-seed experimental boundary.}

\begin{itemize}
  \item A reproducible audio-classification pipeline with fold-aware data
  handling, shared preprocessing, and common model interfaces.
  \item A controlled comparison of baseline and augmented training across three
  architectures and four held-out test conditions.
  \item Deterministic noise corruption with separated training and evaluation
  noise banks and per-sample provenance.
  \item Architecture-level, robustness-curve, and per-class analysis of the
  measured augmentation effects.
\end{itemize}

\section{Related Work}

\DraftingNote{Write this section only after verified BibTeX sources are
available. Cover environmental sound classification, Log-Mel CNN/CRNN models,
ResNet adaptation to spectrograms, audio augmentation and SpecAugment, and
controlled noise-robust evaluation. Compare prior work to this experiment
without claiming an exhaustive systematic review.}

% TODO: Add verified citations for UrbanSound8K, ResNet, GRU, SpecAugment,
% MS-SNSD, and directly relevant environmental-sound robustness studies.

\section{Dataset and Exploratory Analysis}

UrbanSound8K contains 8,732 clips from ten urban sound classes. The official
fold organization was preserved: folds 1--8 were used for training, fold 9 for
validation, and fold 10 for final testing. This prevents a naive random split
from placing closely related source recordings in different subsets.

\begin{table}[t]
  \caption{UrbanSound8K Experimental Split}
  \label{tab:dataset-split}
  \centering
  \begin{tabular}{lrr}
    \toprule
    Subset & Official folds & Samples \\
    \midrule
    Training & 1--8 & 7,079 \\
    Validation & 9 & 816 \\
    Test & 10 & 837 \\
    \midrule
    Total & 1--10 & 8,732 \\
    \bottomrule
  \end{tabular}
\end{table}


\ProjectFigure
  {figures/phase_02/fig_p02_01_class_distribution.png}
  {\columnwidth}
  {UrbanSound8K class-distribution plot}
  {Class distribution of the complete UrbanSound8K inventory. The imbalance
  motivates reporting macro-averaged metrics alongside accuracy.}
  {fig:class-distribution}

\ProjectFigure
  {figures/phase_02/fig_p02_06_representative_log_mel.png}
  {\columnwidth}
  {Representative Log-Mel examples}
  {Representative standardized Log-Mel spectrograms selected
  deterministically from the ten target classes.}
  {fig:representative-log-mel}

\DraftingNote{Add a concise interpretation of duration, sample-rate, and channel
findings from the Phase 2 evidence. Avoid duplicating every EDA plot in the main
paper; move secondary plots to an appendix if required.}

\section{Methodology}

\ProjectFigure
  {figures/phase_03/fig_p03_01_experiment_pipeline.pdf}
  {0.96\columnwidth}
  {Pipeline diagram showing preprocessing, optional training augmentation,
  three model families, and clean/20/10/0 dB evaluation -- pending}
  {Overview of the shared training and controlled robustness-evaluation
  pipeline. Augmentation is applied only to training data.}
  {fig:experiment-pipeline}

\subsection{Audio Preprocessing}

All recordings were converted to mono, resampled to 22,050 Hz, and normalized
to four seconds. Short recordings were zero-padded. Training used seeded random
crops, whereas validation and test processing used deterministic center crops.
The common model input was a per-example standardized Log-Mel spectrogram with
shape $1 \times 64 \times 173$.

\begin{table}[t]
  \caption{Shared Preprocessing Configuration}
  \label{tab:preprocessing}
  \centering
  \begin{tabular}{ll}
    \toprule
    Setting & Value \\
    \midrule
    Sample rate & 22,050 Hz \\
    Clip duration & 4.0 s \\
    Input channels & Mono \\
    FFT/window length & 1,024 \\
    Hop length & 512 \\
    Mel bands & 64 \\
    Dynamic range & 80 dB \\
    Input shape & $1 \times 64 \times 173$ \\
    Normalization & Per-example standardization \\
    \bottomrule
  \end{tabular}
\end{table}


\subsection{Training-Time Augmentation}

The baseline condition disabled robustness-oriented augmentation. The augmented
condition independently applied time shift, random gain, background-noise
mixing, frequency masking, and time masking with probability 0.5. Background
noise used a training-only MS-SNSD noise bank and a uniformly selected SNR in
the configured 0--20 dB range.

\subsection{Model Architectures}

\begin{table}[t]
  \caption{Model Configurations}
  \label{tab:model-parameters}
  \centering
  \begin{tabular}{lrr}
    \toprule
    Architecture & Batch size & Parameters \\
    \midrule
    CNN & 16 & 94,186 \\
    CRNN--BiGRU & 16 & 293,610 \\
    ResNet18 & 8 & 11,175,370 \\
    \bottomrule
  \end{tabular}
\end{table}


Every architecture accepted the same one-channel Log-Mel representation and
returned ten class logits. The ResNet18 model was trained from scratch rather
than initialized with ImageNet weights.

\subsection{Controlled Noise Corruption}

Final evaluation used clean audio and noisy variants at 20, 10, and 0 dB. Noise
was scaled using mean-square signal and noise power:

\begin{equation}
  \mathrm{SNR}_{\mathrm{dB}}
  = 10 \log_{10}\left(\frac{P_{\mathrm{signal}}}{P_{\mathrm{noise}}}\right).
\end{equation}

The corruption seed and sample identifier deterministically selected the noise
file and segment. Therefore, every model received the same corruption for a
given test sample and condition. Training noise and evaluation noise came from
separate MS-SNSD directories.

\ProjectFigure
  {figures/phase_02/fig_p02_08_snr_waveform_comparison.png}
  {\columnwidth}
  {Controlled SNR waveform comparison}
  {Example waveform corruption under the configured clean, 20 dB, 10 dB, and
  0 dB conditions.}
  {fig:snr-waveform-comparison}

\subsection{Metrics}

The study reports accuracy, macro precision, macro recall, macro F1, per-class
precision/recall/F1, and confusion matrices. Robustness summaries include the
clean-to-noisy drop, relative retention, ordinary least-squares slope over the
finite SNR points, and normalized trapezoidal area under the 0--20 dB SNR curve.

\section{Experimental Setup}

All six configurations used deterministic seed 42, CrossEntropyLoss, AdamW
with learning rate 0.001 and weight decay 0.0001, mixed precision on CUDA, a
ReduceLROnPlateau scheduler, and early stopping on validation macro F1.
CNN and CRNN used batch size 16; ResNet18 used batch size 8.

Training ran locally with Python 3.13.1 and PyTorch 2.11.0+cu126 on an NVIDIA
GeForce GTX 1660 Ti. Atomic \texttt{last.pt} checkpoints supported recovery
from electricity interruptions, while final evaluation used only each run's
selected \texttt{best.pt}.

\ProjectFigure
  {figures/phase_03/fig_p03_02_training_validation_curves.png}
  {\columnwidth}
  {Combined loss and validation macro-F1 histories for the six full runs --
  pending}
  {Training and validation histories for the six baseline and augmented
  configurations.}
  {fig:training-curves}

\DraftingNote{Add the exact early-stopping and checkpoint-selection description,
then cite Table~\ref{tab:model-parameters}. Keep pilot and smoke runs out of the
reported experimental results.}

\section{Results}

\begin{table*}[t]
  \caption{Fold-10 Accuracy and Macro F1 Under Controlled Noise}
  \label{tab:complete-test-results}
  \centering
  \begin{tabular}{llrrrrrrrr}
    \toprule
    & & \multicolumn{4}{c}{Accuracy} & \multicolumn{4}{c}{Macro F1} \\
    \cmidrule(lr){3-6}\cmidrule(lr){7-10}
    Architecture & Training & Clean & 20 dB & 10 dB & 0 dB
      & Clean & 20 dB & 10 dB & 0 dB \\
    \midrule
    CNN & Baseline
      & 0.7121 & 0.6703 & 0.6057 & 0.4504
      & 0.7323 & 0.6624 & 0.5857 & 0.4230 \\
    CNN & Augmented
      & 0.7288 & 0.7168 & 0.6989 & 0.5496
      & 0.7469 & 0.7270 & 0.7071 & 0.5481 \\
    CRNN & Baseline
      & 0.7838 & 0.7491 & 0.7085 & 0.4994
      & \textbf{0.8037} & 0.7632 & 0.7134 & 0.4952 \\
    CRNN & Augmented
      & 0.7419 & 0.7491 & 0.7384 & 0.6141
      & 0.7574 & 0.7634 & \textbf{0.7535} & 0.6186 \\
    ResNet18 & Baseline
      & 0.7395 & 0.6977 & 0.5938 & 0.4421
      & 0.7617 & 0.7050 & 0.5774 & 0.4249 \\
    ResNet18 & Augmented
      & 0.7575 & 0.7575 & 0.7264 & 0.6081
      & 0.7781 & \textbf{0.7783} & 0.7432 & \textbf{0.6384} \\
    \bottomrule
  \end{tabular}
\end{table*}

\begin{table*}[t]
  \caption{Macro-F1 Robustness Summary}
  \label{tab:robustness-summary}
  \centering
  \begin{tabular}{llrrrr}
    \toprule
    Architecture & Training & Clean F1 & 0 dB F1 & Normalized SNR AUC & AUC rank \\
    \midrule
    CNN & Baseline & 0.7323 & 0.4230 & 0.5642 & 6 \\
    CNN & Augmented & 0.7469 & 0.5481 & 0.6724 & 3 \\
    CRNN & Baseline & \textbf{0.8037} & 0.4952 & 0.6713 & 4 \\
    CRNN & Augmented & 0.7574 & 0.6186 & 0.7222 & 2 \\
    ResNet18 & Baseline & 0.7617 & 0.4249 & 0.5712 & 5 \\
    ResNet18 & Augmented & 0.7781 & \textbf{0.6384} & \textbf{0.7258} & 1 \\
    \bottomrule
  \end{tabular}
\end{table*}

\begin{table*}[t]
  \caption{Augmented-Minus-Baseline Effects}
  \label{tab:augmentation-effects}
  \centering
  \begin{tabular}{lrrrrr}
    \toprule
    Architecture & Clean accuracy & Clean macro F1
      & 0 dB accuracy & 0 dB macro F1 & Macro-F1 AUC \\
    \midrule
    CNN & +0.0167 & +0.0146 & +0.0992 & +0.1252 & +0.1082 \\
    CRNN & -0.0418 & -0.0464 & +0.1147 & +0.1234 & +0.0509 \\
    ResNet18 & +0.0179 & +0.0163 & +0.1661 & +0.2135 & +0.1546 \\
    \bottomrule
  \end{tabular}
\end{table*}


\ProjectFigure
  {figures/phase_03/fig_p03_03_macro_f1_robustness_curves.png}
  {\columnwidth}
  {Macro-F1 robustness curves}
  {Macro-F1 performance under clean and controlled background-noise
  conditions. Solid lines denote augmented training and dashed lines denote
  baseline training.}
  {fig:macro-f1-robustness}

\ProjectFigure
  {figures/phase_03/fig_p03_05_macro_f1_augmentation_effects.png}
  {\columnwidth}
  {Augmentation-effect comparison}
  {Architecture-matched macro-F1 difference between augmented and baseline
  training at each evaluation condition.}
  {fig:augmentation-effects}

\ProjectFigure
  {figures/phase_03/fig_p03_06_macro_f1_robustness_auc.png}
  {\columnwidth}
  {Robustness-AUC comparison}
  {Normalized macro-F1 SNR area under the noisy-condition curve. Higher values
  indicate stronger performance across 0, 10, and 20 dB.}
  {fig:robustness-auc}

CRNN baseline obtained the highest clean macro F1 (0.8037). ResNet18 augmented
obtained the highest normalized macro-F1 SNR AUC (0.7258) and the highest 0 dB
macro F1 (0.6384). The three augmented configurations occupied the top three
robustness-AUC ranks.

\ProjectFigure
  {figures/phase_03/fig_p03_07_zero_db_per_class_effects.png}
  {\columnwidth}
  {Zero-dB class-level effects}
  {Per-class F1 change from augmentation at 0 dB for each architecture.}
  {fig:zero-db-class-effects}

\ProjectFigure
  {figures/phase_03/fig_p03_08_confusion_matrix_comparison.png}
  {\columnwidth}
  {Normalized clean and 0 dB confusion-matrix comparison -- pending}
  {Normalized confusion matrices for the selected baseline and augmented
  configurations under clean and 0 dB conditions.}
  {fig:confusion-comparison}

\DraftingNote{Expand this section by describing the tables and figures without
repeating every cell. Do not use statistically significant language until
confidence intervals or tests exist.}

\section{Discussion}

The completed experiment supports the hypothesis that training-time
augmentation can reduce degradation under controlled test noise. The magnitude
was architecture-dependent. ResNet18 showed the largest normalized macro-F1
SNR-AUC gain (0.1546), while CRNN showed a clean-versus-robustness trade-off:
its clean macro F1 decreased by 0.0464, but its 0 dB macro F1 increased by
0.1234.

\DraftingNote{Interpret why model capacity and temporal modeling may produce
different augmentation responses, but label mechanisms as hypotheses rather
than proven causes. Discuss the strong ResNet18 0 dB class gains and the
remaining difficult classes. Relate findings to verified prior work only after
the bibliography is populated.}

\section{Limitations and Threats to Validity}

\begin{itemize}
  \item Final testing used one official UrbanSound8K fold rather than complete
  ten-fold cross-validation.
  \item Each of the six training configurations currently has one completed
  training seed.
  \item Confidence intervals and formal paired statistical tests have not yet
  been reported.
  \item Evaluation used one held-out noise collection containing 51 files and
  three finite noisy SNR values.
  \item ResNet18 was non-pretrained, and modern pretrained audio transformers
  were outside the experiment scope.
  \item Conclusions are limited to this controlled split and corruption
  protocol and should not be generalized to every acoustic environment.
\end{itemize}

\section{Conclusion}

\DraftingNote{Write the conclusion after Results and Discussion are finalized.
Answer both research questions directly: augmentation improved noisy-condition
robustness for all three architectures, but its magnitude and clean-data
trade-off differed. Mention ResNet18 augmented as the strongest overall
robustness model and propose repeated seeds, confidence intervals, broader noise
sets, and fold-based evaluation as future work.}


\bibliographystyle{IEEEtran}
\bibliography{bibliography/references}

\end{document}
